
\section{The Neural TKF Model}
\label{sec:neural-tkf}

\subsection{Beyond HMMs: epistasis and neural networks}

A fundamental limitation of Hidden Markov Model (HMM) approaches to sequence evolution
is their inability to model \emph{epistasis}---the phenomenon whereby
the evolutionary constraints at one site depend on the identity of residues at other sites.
In an HMM, emission probabilities at each position are conditionally independent
given the hidden state, so the substitution process at site $i$
cannot depend on the residue at site $j$.
This is a poor approximation for proteins (where stability and function
depend on networks of interacting residues)
and for structured RNAs (where base-pairing imposes strong covariation between paired sites).

Stochastic context-free grammars (SCFGs) offer a partial remedy:
by emitting paired terminals, they can capture the covariation
between base-paired positions in RNA secondary structure.
However, SCFGs impose a highly constrained, tree-structured form of epistasis:
only positions that are co-emitted by the same nonterminal can covary,
and the pattern of dependencies is rigidly determined by the grammar's derivation tree.
SCFGs cannot model the dense, long-range epistatic networks
found in protein tertiary structure or RNA pseudoknots.

The strongest candidate for moving beyond these limitations
is \emph{neural network parameterization} of the evolutionary process.
By conditioning transition and emission parameters on learned embeddings
of the full sequence context, neural models can in principle capture
arbitrary epistatic interactions, limited only by the capacity of the network.
The challenge is to do so while retaining the interpretable probabilistic structure---alignment,
phylogeny, and time-reversibility---that makes statistical models of evolution scientifically useful.

\subsection{Neural transducers}

A \emph{neural transducer} replaces the fixed parameters of a statistical sequence transducer
(such as a Pair HMM) with position-specific parameters predicted by a neural network.
Given an input sequence $X = (x_1, \ldots, x_L)$ and evolutionary time $\evoltime$,
a neural embedding network $\mathcal{E}(X)$ produces per-site embeddings,
from which prediction heads derive site-specific transition and emission parameters.
The model likelihood is computed by a standard Forward algorithm
over the transducer's state space,
but with neural-predicted parameters substituted at each position.
Training proceeds by differentiating the Forward log-likelihood with respect to the network weights.

This approach occupies a middle ground between two extremes.
At one end are classical homogeneous models (TKF91, TKF92, and their mixtures)
with shared parameters across all sites.
At the other end are unconstrained sequence-to-sequence models
that predict descendant tokens autoregressively
without any explicit alignment or evolutionary structure~\cite{KoehlEtAl2026PEINT}.
The neural transducer retains the alignment structure and time-reversibility
of the underlying stochastic process,
while allowing site-specific evolutionary rates
that can capture context-dependent selection pressures.

Viewed through this lens,
a neural transducer with a $k$th-order Markov embedding
can be seen as a $(k+1)$th-order approximation
to a whole-sequence evolutionary model.
Models such as Cosine~\cite{LuEtAl2026Cosine}
also condition on the full ancestor sequence via attention,
but use the alignment purely as an attention guide
rather than as an explicit latent variable in the generative model.
Models such as PEINT~\cite{KoehlEtAl2026PEINT}
dispense with alignment entirely,
gaining flexibility at the cost of interpretable evolutionary structure.
The neural transducer combines the alignment-guided architecture
of Cosine with the epistatic sequence-to-sequence flexibility of PEINT,
while retaining a well-defined probabilistic indel model
that can be composed along a phylogenetic tree.

\subsection{The Neural TKF92 model}

Large and Holmes~\cite{LargeHolmes2026} introduced the Neural TKF model,
which parameterizes a site-specific TKF92 process using neural network embeddings.
Given an ancestor--descendant pair $(X, Y)$ separated by time $\evoltime$,
the model works as follows.

\paragraph{Sequence embeddings.}
Separate embedding networks $\mathcal{E}_{\mathrm{anc}}(X)$
and $\mathcal{E}_{\mathrm{desc}}(Y)$
produce per-site embedding vectors.
The ancestor embedder uses bidirectional context
(the full ancestor is observed),
while the descendant embedder uses causal (left-to-right) context
to maintain a valid autoregressive factorization.
Several architectures were explored, including residual CNNs,
bidirectional LSTMs, and Transformers with rotary positional encoding;
the Transformer architecture performed best.

\paragraph{Prediction heads.}
At each aligned position $(i, j)$,
the ancestor and descendant embeddings are concatenated
along with the previous alignment state and evolutionary time $\evoltime$,
then passed through feedforward networks to predict site-specific TKF92 parameters:
\begin{itemize}[nosep]
\item Equilibrium distribution $\eqm_{ij}$ (via softmax over alphabet logits);
\item Deletion rate $\delrate_{ij}$ (via a bounded sigmoid $\sigma_{\mathrm{bound}}$);
\item Insertion--deletion offset $\delta_{ij}$ with $\insrate_{ij} = \delrate_{ij}(1 - \delta_{ij})$,
  ensuring $\delrate > \insrate$ (also via bounded sigmoid);
\item Fragment extension probability $\ext_{ij}$ (via standard sigmoid).
\end{itemize}
The bounded sigmoid $\sigma_{\mathrm{bound}}(x, b_{\min}, b_{\max})
= b_{\min} + (b_{\max} - b_{\min}) / (1 + e^{-x})$
prevents rates from diverging during training.

\paragraph{Emission and transition.}
Emissions follow a site-specific F81 substitution model:
for a match column $(a, b)$ with $a, b \in \alphabet$,
the emission probability is $[\exp(\exch_{ij} \evoltime)]_{ab}$,
where $\exch_{ij}$ is the F81 rate matrix constructed from $\eqm_{ij}$.
Insert emissions use $(\eqm_{ij})_b$.
The transition matrix is the standard TKF92 matrix
(with states $\{ \mat, \ins, \del, \fin \}$),
but with site-specific parameters $(\insrate_{ij}, \delrate_{ij}, \ext_{ij})$
substituted at each position.

\paragraph{Training.}
The pairwise alignment between ancestor and descendant
serves as the latent path through the Pair HMM.
The Forward algorithm computes the marginal likelihood
$P(Y | X, \evoltime)$ by summing over all valid alignments,
and the network is trained by gradient descent on $-\log P(Y | X, \evoltime)$.
The Forward dynamic program is fully differentiable.

\paragraph{Alignment-guided attention.}
A key architectural feature is that the pairwise alignment
can serve as a \emph{causal attention guide}:
at descendant position $j$,
cross-attention to the ancestor is restricted
to positions $i' \leq i$ where $i$ is the rightmost ancestor position
aligned to or before $j$.
This alignment-guided masking provides an inductive bias
that encourages the network to attend to evolutionarily related positions,
while maintaining the causal structure needed for autoregressive decoding.

\paragraph{Extensions.}
The embedding networks can attend to richer inputs
beyond the raw sequence:
multiple sequence alignments (MSAs) of homologs,
predicted or experimentally determined 3D structures,
or functional annotations.
These additional inputs are concatenated or cross-attended
into the per-site embeddings before the prediction heads,
allowing the model to condition evolutionary rates
on structural and functional context.

\subsection{The Triad model for ancestral reconstruction}
\label{sec:triad}

The Neural TKF model as described above is trained on ancestor--descendant pairs,
where both sequences are observed (e.g., extant homologs related by a known branch).
For ancestral reconstruction on a phylogenetic tree,
we require the reverse direction:
inferring an ancestor given its descendants.

The \emph{Triad model} addresses this by generating an ancestral sequence
conditioned on two sibling sequences.
Specifically, given a tree of the form
\[
\texttt{((Sib}_1\texttt{:}t_1\texttt{, Sib}_2\texttt{:}t_2\texttt{)Parent:}t_3\texttt{)Grandparent;},\]
the model maximizes $P(\text{Grandparent} \mid \text{Sib}_1, \text{Sib}_2, t_1, t_2, t_3)$,
using the grandparent as a proxy for the unobserved parent
(in the limit $t_3 \to 0$, the grandparent approaches the parent).
Three branch lengths $(t_1, t_2, t_3)$ replace the single time parameter,
extending the dimensionality of the evolutionary conditioning.

\paragraph{Input representation.}
The two sibling sequences are first aligned to each other,
producing a pairwise alignment over a $(|\alphabet|+1)^2 - 1$ token alphabet
(every pair of aligned characters including gaps, excluding the double-gap).
This sibling alignment is treated as a fixed input tape:
each column is an atomic symbol,
and the semantics of it being a pairwise alignment are opaque to the decoder.

\paragraph{Output representation.}
The output alphabet is the same alignment-augmented alphabet
as in the pairwise Neural TKF:
$|\alphabet|$ match tokens, $|\alphabet|$ insert tokens,
a delete token, and an end-of-sequence token.
Here ``insert'' and ``delete'' refer to the grandparent tape
relative to the sibling tape;
the reversal of evolutionary direction
(siblings$\to$grandparent vs.\ ancestor$\to$descendant)
is handled naturally by the model's time-reversibility.

\paragraph{Full alignment map.}
Unlike the pairwise model, where the alignment between ancestor and descendant
is a Markov chain of $\mat/\ins/\del$ states,
the Triad model maintains a \emph{full alignment map}
between the grandparent output and the sibling input tape.
This alignment map records, for each grandparent position,
which sibling-tape columns it aligns to.
The full map (rather than just the most recent alignment state)
is made available to the decoder,
allowing attention to all previously aligned positions.

In the decoder, this alignment map augments both:
\begin{itemize}[nosep]
\item \emph{Cross-attention}: at each grandparent output position,
  the model attends to sibling-tape positions
  guided by the cumulative alignment,
  giving access to covariant positions
  (e.g., compensatory base-pair partners in RNA);
\item \emph{Self-attention}: the alignment map can also modulate
  self-attention within the grandparent sequence,
  so that positions aligned to the same or nearby sibling-tape columns
  are encouraged to attend to each other.
\end{itemize}

This full-map attention is a strictly richer conditioning
than the Markovian alignment state used in the pairwise model.
It enables the Triad decoder to exploit long-range covariation
in the sibling sequences when reconstructing the ancestor,
which is particularly important for structured RNAs
where compensatory mutations at paired positions carry strong phylogenetic signal.

\paragraph{Training data.}
Training triads $(\text{Grandparent}, \text{Sib}_1, \text{Sib}_2)$
are sampled from phylogenetic trees with known alignments
(e.g., Rfam seed alignments for RNA families).
The sampling targets uniform coverage over sibling divergences
$(0 < t_1 \leq 1, 0 < t_2 \leq 1)$
while keeping the parent--grandparent distance $t_3$ short,
since at inference time the model operates in the limit $t_3 \to 0$
for progressive ancestral reconstruction.

\paragraph{Chapman--Kolmogorov consistency.}
A key test for any time-dependent evolutionary model
is Markovian consistency:
evolving from time 0 to $t_1 + t_2$ directly
should give the same distribution as evolving from 0 to $t_1$
and then from $t_1$ to $t_1 + t_2$.
For the Neural TKF model, this can be tested
by using a TKF92-based importance sampler
to estimate the marginal distribution
$P(Y | X, t_1 + t_2) \approx \sum_Z P(Z | X, t_1) P(Y | Z, t_2)$,
where $Z$ ranges over intermediate sequences
sampled from the first-leg conditional.
Discrepancies from the direct model $P(Y | X, t_1 + t_2)$
reveal failures of Markovian consistency,
which are expected for finite-order neural approximations
and can guide architecture improvements.
